% P01 — A Trust Calculus for Attestation Tiers — arXiv preprint
% Build: tectonic -Z shell-escape main.tex  (XeTeX engine; fetches packages)
% STATUS: working draft. The results sections read generated LaTeX from
% ../results/. Do not type numbers into this file.
%
% Preamble, colour palette, \aaimark and the cover-page TikZ block are copied
% from ../../ov-poc-standard/paper/main.tex so the two papers look like one
% series. One deliberate difference: that paper loads `minted`, which needs a
% pygmentize binary on the build host. This paper has no syntax-highlighted
% listings, so minted is not loaded and the build has one fewer external
% dependency. `listings` is configured identically in case that changes.
\documentclass[11pt]{article}

\usepackage[margin=1in]{geometry}
\usepackage{tikz}
\usepackage{graphicx}
\usepackage{booktabs}
\usepackage{array}
\usepackage{amsmath, amssymb}
\usepackage{listings}
\definecolor{codepanel}{HTML}{F7F5F1}
\usepackage{xcolor}
\usepackage{colortbl}
\usepackage[hidelinks]{hyperref}
\usepackage{enumitem}
\usepackage{microtype}
\usepackage{caption}
\usepackage{siunitx}
\usepackage{amsthm}
\newtheorem{definition}{Definition}
\newtheorem{theorem}{Theorem}
\newtheorem{proposition}{Proposition}
\usepackage[section]{placeins}

% Advanced AI Society brand palette
\definecolor{aailime}{HTML}{CFFF04}
\definecolor{aailimetint}{HTML}{F2FFB8}
\definecolor{aaiink}{HTML}{0A0A0A}
\definecolor{aaired}{HTML}{FFE3E4}
\definecolor{aaineutral}{HTML}{E7E3DD}
\definecolor{aaipaper}{HTML}{F0EDEA}
\definecolor{aaimuted}{HTML}{8F8A82}
\definecolor{aaidarklime}{HTML}{7A9900}

\lstset{basicstyle=\ttfamily\footnotesize, breaklines=true, frame=single,
        framerule=0.3pt, xleftmargin=1em, columns=fullflexible}
\captionsetup{font=small, labelfont=bf}

\newcommand{\parallax}{\texttt{parallax}}
\newcommand{\Never}{\infty}
% An underscore inside a code identifier, with a break opportunity after it.
% Capability names run to 32 characters and will not fit a table column
% otherwise; TeX will not hyphenate inside \texttt, so the break points have
% to be stated.
\newcommand{\ub}{\_\allowbreak}

% Counts the prose states, generated from the artifact rather than typed.
%
% Every other generated file here is a table body. These are `\newcommand`s,
% because generating tables only stops a *table* drifting from the code -- it
% does nothing about the sentence beside the table stating the same number in
% words. Twice now that is exactly what went wrong. The abstract, Section 1
% and the Conclusion all said three of the five parties behind a TDX
% measurement are undetectable, while three generated tables on the facing
% pages said four. And Section 6, in the paragraph written to concede what
% the paper is not entitled to claim, said Sigma_1 declares the fewest
% mechanisms when it declares the most. Both families of count now come from
% macros and cannot be typed wrong.
%
% What this does not fix: Section 1 also names the undetectable parties one
% by one. If the count changed, the macro would update the number and leave
% the list of names stale, so `the_tdx_headline_count_is_four` in src/tex.rs
% fails on any such change to force a human to look. Likewise
% `one_assumption_per_mechanism_would_invert_the_cardinality_ranking` in
% src/tiers.rs pins the *shape* of Section 6's concession, which no macro
% covers.
% Generated by `parallax solve --format latex-counts` via scripts/regen-results.sh — do not edit.
% Macro definitions for the paper's prose, \input from its preamble.
\newcommand{\SigmaOneAssumptions}{three}
\newcommand{\SigmaOneParties}{three}
\newcommand{\SigmaOneUndetectable}{three}
\newcommand{\SigmaOneBounded}{zero}
\newcommand{\SigmaOneMechanisms}{three}
\newcommand{\SigmaOneMechanismKinds}{two}
% Not every macro above is used today. They are the counts a sentence
% about a deployment tends to reach for, and defining an unused one
% costs nothing next to typing a used one by hand.

% Generated by `parallax solve --format latex-counts` via scripts/regen-results.sh — do not edit.
% Macro definitions for the paper's prose, \input from its preamble.
\newcommand{\SigmaTwoAssumptions}{five}
\newcommand{\SigmaTwoParties}{five}
\newcommand{\SigmaTwoUndetectable}{four}
\newcommand{\SigmaTwoBounded}{one}
\newcommand{\SigmaTwoMechanisms}{one}
\newcommand{\SigmaTwoMechanismKinds}{one}
% Not every macro above is used today. They are the counts a sentence
% about a deployment tends to reach for, and defining an unused one
% costs nothing next to typing a used one by hand.

% Generated by `parallax solve --format latex-counts` via scripts/regen-results.sh — do not edit.
% Macro definitions for the paper's prose, \input from its preamble.
\newcommand{\SigmaThreeAssumptions}{eleven}
\newcommand{\SigmaThreeParties}{eleven}
\newcommand{\SigmaThreeUndetectable}{seven}
\newcommand{\SigmaThreeBounded}{four}
\newcommand{\SigmaThreeMechanisms}{three}
% Not every macro above is used today. They are the counts a sentence
% about a deployment tends to reach for, and defining an unused one
% costs nothing next to typing a used one by hand.

% Generated by `parallax solve --format latex-counts` via scripts/regen-results.sh — do not edit.
% Macro definitions for the paper's prose, \input from its preamble.
\newcommand{\SigmaFourAssumptions}{five}
\newcommand{\SigmaFourParties}{five}
\newcommand{\SigmaFourUndetectable}{three}
\newcommand{\SigmaFourBounded}{two}
\newcommand{\SigmaFourMechanisms}{two}
% Not every macro above is used today. They are the counts a sentence
% about a deployment tends to reach for, and defining an unused one
% costs nothing next to typing a used one by hand.

% Generated by `parallax solve --format latex-counts` via scripts/regen-results.sh — do not edit.
% Macro definitions for the paper's prose, \input from its preamble.
\newcommand{\SigmaFiveAssumptions}{eight}
\newcommand{\SigmaFiveParties}{seven}
\newcommand{\SigmaFiveUndetectable}{seven}
\newcommand{\SigmaFiveBounded}{one}
\newcommand{\SigmaFiveMechanisms}{two}
% Not every macro above is used today. They are the counts a sentence
% about a deployment tends to reach for, and defining an unused one
% costs nothing next to typing a used one by hand.


% Placeholders, visible in the PDF while \ifdraft is true.
%
% A LaTeX comment marking unfinished content is invisible to everyone reading
% the compiled paper, which is everyone who reviews it. These render as a
% marked block instead, so a placeholder cannot be circulated unnoticed.
% Setting \draftfalse before submission turns them off and, more usefully,
% makes any surviving placeholder disappear from the text rather than
% silently becoming a claim -- so grep for \placeholder before submitting.
\newif\ifdraft
\drafttrue
\newcommand{\placeholder}[1]{%
  \ifdraft
    \par\vspace{0.4em}\noindent
    \fcolorbox{aaidarklime}{aailimetint}{%
      \parbox{\dimexpr\linewidth-2\fboxsep-2\fboxrule\relax}{%
        \small\sffamily\raggedright
        \textbf{Placeholder.} #1}}%
    \par\vspace{0.4em}
  \fi}

% The Society's bracket mark, drawn natively (lime on the dark band)
\newcommand{\aaimark}[1]{%
\begin{tikzpicture}[x=#1, y=#1, line cap=butt]
  \fill[aailime] (0,0) rectangle (0.16,1);
  \fill[aailime] (0.16,0) rectangle (0.42,0.24);
  \fill[aailime] (0.16,0.76) rectangle (0.42,1);
  \fill[aailime] (1.04,0) rectangle (1.20,1);
  \fill[aailime] (0.78,0) rectangle (1.04,0.24);
  \fill[aailime] (0.78,0.76) rectangle (1.04,1);
\end{tikzpicture}}

\begin{document}

% ---------------------------------------------------------------- cover page
\thispagestyle{empty}

\begin{tikzpicture}[remember picture, overlay]
  \fill[aaiink] (current page.north west) rectangle
    ([yshift=-6.35cm]current page.north east);
  \fill[aailime] ([yshift=-6.35cm]current page.north west) rectangle
    ([yshift=-6.53cm]current page.north east);
  \node[anchor=north west, inner sep=0pt]
    at ([xshift=1in, yshift=-0.95cm]current page.north west) {\aaimark{0.62cm}};
  \node[anchor=north west, inner sep=0pt, text width=14.2cm]
    at ([xshift=1in+1.15cm, yshift=-0.98cm]current page.north west)
    {\sffamily\bfseries\footnotesize\color{aailime}%
     ADVANCED AI SOCIETY\; \textcolor{aaimuted}{$\cdot$}\; PROOF-OF-CONTROL
     INITIATIVE\\[3pt]
     {\normalfont\sffamily\footnotesize\color{aaimuted}
      Working draft v0.1 $\cdot$ August 2026 $\cdot$ for co-author review
      prior to submission}};
  \node[anchor=north west, inner sep=0pt, text width=\textwidth]
    at ([xshift=1in, yshift=-2.55cm]current page.north west)
    {\sffamily\bfseries\fontsize{21.5}{26}\selectfont\color{aaipaper}%
     A Trust Calculus for Attestation Tiers: Computing Residual Trust and
     Testing Whether Deployments Can Be Ranked};
  \node[anchor=north west, inner sep=0pt, text width=\textwidth]
    at ([xshift=1in, yshift=-5.55cm]current page.north west)
    {\sffamily\small\color{aaipaper}%
     \textcolor{aailime}{\rule[0.5ex]{1.6em}{2pt}}\hspace{0.6em}%
     Which parties must you still trust, and for how long would you not know?};
\end{tikzpicture}

\vspace*{5.9cm}

% Compact author block
{\small\noindent
\textbf{Jim Schwoebel}\textsuperscript{1,2}

\smallskip\noindent
{\footnotesize
\textsuperscript{1}\,Quome \quad
\textsuperscript{2}\,Advanced AI Society, \url{https://advancedaisociety.org}%
\footnote{Prepared within the Proof-of-Control Initiative of the Advanced AI
Society. Authorship is limited to contributors to this manuscript; additional
co-authors will be added only upon explicit consent and substantive
contribution.}}}

\vspace{0.9em}
\noindent\textcolor{aailime}{\rule{\textwidth}{1.6pt}}
\vspace{0.2em}

\begin{center}
{\sffamily\bfseries\small ABSTRACT}
\end{center}
\vspace{-0.5em}
\noindent{\small
Attestation systems are routinely described as removing the need to trust
anybody. The Proof-of-Control standard~\cite{pocspec} says of its Tier~3 that
``there is nobody left to trust; anyone can check the evidence with published
tools.''
That sentence is false, and this paper shows it is false rather than
suspecting it. Deploying on Intel TDX and enumerating what a verifier must
accept to believe one measurement yields \SigmaTwoParties{} distinct parties,
and \SigmaTwoUndetectable{} of the \SigmaTwoParties{} have no detection
mechanism at all: if they misbehave, nothing in
the system will ever say so. The fifth is bounded only by how often
revocation data is refreshed. We give a small calculus that computes this set
instead of arguing about it. A judgment $\Gamma \vdash P : T, \Delta$ reads
``proposition $P$ holds provided the parties in $T$ are honest, and a
violation would be detected within $\Delta$''; mechanisms compose by joining
$T$ under union and $\Delta$ under max, the former over a trust-set lattice
of finite height. We implement it in \parallax{}, a Rust tool that reads a
deployment description and prints the residual trust set, and we show that
termination on cyclic delegation --- the case the original design proposed to
handle with tabled resolution --- follows from the finite height of that one
lattice and needs no tabling. Of the deployments we encoded, the two pairs
that attest the same proposition and can therefore be compared at all come
out \emph{incomparable} under set inclusion rather than merely different,
which would mean no ordinal tier ladder can rank them. We then run each of the
four ways we can think of to order deployments by verifiability --- our own
list, not a survey --- against those same
encodings. All four fail, and each fails differently: ordering by the size of
the trust set ranks a bare software host above a witness quorum, which inverts
the word. We are explicit about the
limit that matters most: the solver computes over the dependencies somebody
wrote down, and cannot discover one that was left out.}

\vspace{0.6em}
\noindent{\small\textbf{Artifact.} \parallax{}~\cite{parallaxtool} is
Apache-2.0. Every \emph{table and figure} in
Sections~\ref{sec:eval} and~\ref{sec:tiers} is regenerated from the
tool by \texttt{scripts/\allowbreak regen-results.sh} and read into this document from
the generated files; continuous integration fails if a committed result
disagrees with the code that produced it. Sentences near those tables
sometimes restate a figure from one, and those sentences are typed rather
than generated; the ones that could drift are pinned by tests named in
\texttt{paper/\allowbreak README.md}. Figures elsewhere in the paper ---
deployment parameters, code size, the size of the termination example ---
are stated in the prose and are not tool output.}

\clearpage

\tableofcontents
\clearpage

% ================================================================
\section{The Problem}
\label{sec:intro}

Somebody hands you a signed attestation. It says a particular piece of
software, with a particular measurement, ran inside a hardware-protected
environment on a machine you do not own~\cite{intelTdx}. The pitch is that you
no longer have
to take anyone's word for it: the hardware vouches for the software, the
signature vouches for the hardware, and you can check the whole thing yourself
with published tools.

So check it. Here is what you actually did.

You verified a signature chain rooted in a key that Intel provisioned. You
believed that key was issued to real silicon, because Intel's certification
service said so. You believed the quote was produced by a quoting enclave
that measured the domain honestly and did not leak its signing key. You
believed the host machine did not inject a different firmware measurement at
boot. And you believed that the reference value you compared against ---
the number that says ``this measurement means \emph{that} code'' --- was
published by somebody who knows what code it belongs to.

That is \SigmaTwoParties{} parties. You did not remove trust; you moved it. And the thing
that should stop the conversation is not the count but the fourth column of
Table~\ref{tab:tdx}: for \SigmaTwoUndetectable{} of those
\SigmaTwoParties{}, there is no detection mechanism
at all. If Intel's silicon is subverted, or the quoting enclave's key leaks,
or the host injects a firmware measurement of its own choosing, or the
reference values are wrong, the attestation still verifies
cryptographically. Nothing anywhere in the system contradicts it. The failure
is silent and it is permanent. The only one of the five you would ever hear
about is Intel's certification service, and only because its collateral goes
stale on a schedule.

\begin{table}[h]
\centering\small
\caption{What a verifier must accept to believe a single Intel TDX
measurement, derived by deploying the Proof-of-Control reference
implementation and enumerating what was left. The fourth column is the point.}
\label{tab:tdx}
\setlength{\tabcolsep}{4pt}
\begin{tabular}{@{}>{\raggedright\arraybackslash}p{2.4cm}
                  >{\raggedright\arraybackslash}p{4.2cm}
                  >{\raggedright\arraybackslash}p{4.8cm}
                  >{\raggedright\arraybackslash}p{2.9cm}@{}}
\toprule
\textbf{Party} & \textbf{What you must assume} &
\textbf{What a breach buys the adversary} & \textbf{Detectable in} \\
\midrule
Intel &
the TDX module, microcode and CPU behave as specified; the provisioning key
is not misissued &
arbitrary quote forgery, state extraction, memory manipulation inside the
hardware boundary &
\textcolor{aaidarklime}{\textbf{never}} \\
\addlinespace
Intel's certification service (PCS) &
the collateral fetched to check a quote is authentic (TCB recovery info,
CRLs, QE certs) &
revoked, vulnerable or emulated hardware presented as a valid target &
collateral refresh interval ($\approx$\,\SI{12}{\hour}) \\
\addlinespace
The quoting enclave &
it measures TD reports accurately and signs without key leakage &
forged execution claims on a compromised host &
\textcolor{aaidarklime}{\textbf{never}} \\
\addlinespace
The cloud operator &
the host does not extract domain memory outside the TDX threat model; the
firmware measured into RTMRs is what it claims &
deceptive measurement injection at boot; side-channel extraction &
\textcolor{aaidarklime}{\textbf{never}} \\
\addlinespace
Whoever publishes reference values &
the MRTD compared against belongs to the code you believe it does &
valid attestation of subverted software inside a genuine TD &
\textcolor{aaidarklime}{\textbf{never}} \\
\bottomrule
\end{tabular}
\end{table}

Anchoring the evidence log to a public ledger removes none of these. It is a
real and valuable property --- it makes the \emph{history} publicly auditable
--- but the \emph{measurement} still rests on the chain above. And, as
Section~\ref{sec:formal} shows, anchoring does not even come for free: it
ingests a new assumption about the settlement layer's execution fidelity
while it is converting the old one.

\paragraph{Two properties that are being conflated.} The tier ladder runs
together two things a careful reader will separate:

\begin{itemize}[itemsep=1pt,topsep=2pt]
\item \textbf{Public verifiability} --- anyone may check the evidence, with
published tools, without permission from the party being checked.
\item \textbf{Trust independence} --- no named party's honesty is
load-bearing for the conclusion.
\end{itemize}

Hardware attestation buys the first outright. It leaves the second largely
intact. Confusing them is what produces a sentence like ``there is nobody
left to trust,'' and that sentence sits at the binary threshold the whole
standard is organized around. If the threshold is mis-drawn, every
conformance claim above it says less than its definition promises.

\paragraph{Arguing about this is not working.} Everyone in the field agrees,
when pressed, that hardware roots are assumptions. What nobody has is a way to
settle a specific disagreement about a specific deployment. Reviewer says
``you are still trusting the build pipeline''; vendor says ``no, the
attestation covers that''; and there the matter rests, because there is no
procedure that takes a system description and returns an answer. That is what
we set out to build.

\paragraph{What we did.} We wrote down a judgment form, gave composition rules
for each mechanism the standard admits, implemented them, and ran them on
five encoded deployments.

\begin{enumerate}[itemsep=1pt,topsep=2pt]
\item \textbf{The calculus} (Section~\ref{sec:formal}). A judgment $\Gamma
\vdash P : T, \Delta$: the proposition $P$ holds provided the parties in $T$
are honest, with any violation detectable within $\Delta$. Composition is
join in two lattices --- union on trust sets, max on detection latencies.
\item \textbf{The tool} (Section~\ref{sec:formal}, Section~\ref{sec:eval}).
\parallax{}, about 1\,800 lines of Rust excluding tests and roughly twice
that again in tests, reads a deployment as TOML and
emits the residual trust set, a machine-readable manifest, a comparison
between two deployments, a shared-dependency analysis, and a policy check.
\item \textbf{Termination, the easy way} (Section~\ref{sec:termination}). The
original design proposed tabled (SLG) resolution so that circular delegation
would terminate. It is not needed. The fixpoint lives in the powerset of a
finite assumption set, which has finite height, so bottom-up evaluation halts
regardless of cycles. We say why the harder machinery solves a problem we do
not have.
\item \textbf{What we found} (Section~\ref{sec:eval}). The tool reproduces
Table~\ref{tab:tdx} --- each party's capability, detection bound and impact,
and the fact that exactly these five are load-bearing --- from a deployment
description that says none of it, and it finds a shared build-pipeline
dependency underneath a hybrid TEE-plus-ZK design that is marketed as defence
in depth. Of the encoded deployments, the two pairs that can be compared at
all --- comparison requires both sides to attest the same proposition ---
come out \emph{incomparable}.
\item \textbf{Whether the ladder can be repaired} (Section~\ref{sec:tiers}).
Four candidate orderings, run against the four non-hybrid
deployments. All four fail,
each in its own way: counting assumptions ranks the systems backwards, set
inclusion ranks almost nothing, composed detection latency is $\Never$ for
every deployment and so ranks everything the same, and collusion cost is not
derivable from a deployment description at all.
\item \textbf{What is wrong with all of this} (Section~\ref{sec:open}). Led by
the completeness problem, which we think is the strongest objection to the
whole approach and which we have not solved.
\end{enumerate}

\paragraph{One thing we are not claiming.} \parallax{} has been run on five
deployments, all of which we wrote. It reproduces a trust set we had already
derived by hand, and it found a shared dependency in a deployment we
constructed. It has not been pointed at somebody else's production system,
and no claim in this paper should be read as saying it has. Section
\ref{sec:eval} says exactly which examples exist and Section~\ref{sec:open}
says what that leaves open.

% ================================================================
\section{What Others Have Built}
\label{sec:related}

\paragraph{The roles exist; the obligations do not.} The IETF's remote
attestation architecture, RFC 9334~\cite{rfc9334}, already gives us the
vocabulary: Attester, Verifier, Relying Party, Endorser, Reference Value
Provider. This is the closest existing thing and we adopt its names. But
RFC~9334 is \emph{descriptive}. It says who plays what part in an attestation
exchange. It does not say what breaks when a part is played dishonestly, and
that --- not the role taxonomy --- is what a relying party needs. Reading
RFC~9334 tells you an Endorser exists. It does not tell you that if the
Endorser lies, nothing anywhere in the protocol will ever notice.

\paragraph{Authorization logics: the right primitives, the wrong direction.}
Abadi, Burrows, Lampson and Plotkin's calculus for access
control~\cite{abadi1993}, and the Taos work that grew from
it~\cite{lampson1992}, gave us \textsf{says} and \textsf{speaks-for}. Those
are exactly the right shape for ``Intel says this measurement,'' and our
delegation relation is \textsf{speaks-for} under another name. Nexus
authorization logic~\cite{nal} is the closest prior art of all: an
authorization logic built specifically for attestation-based systems. NAL
divides the bases for believing a claim into three kinds --- \emph{axiomatic}
(believed by policy fiat), \emph{analytic} (derived by mechanical analysis of
code), and \emph{synthetic} (produced by a trusted transformation or an
isolated execution environment). That is already a taxonomy of \emph{why} a
claim is believed, and it maps onto what we compute: every assumption in a
residual trust set is, in NAL's terms, an axiomatic or synthetic basis that
survived. What NAL does not do is emit the set.

Nor do the declarative policy languages. SecPAL~\cite{secpal},
DKAL~\cite{dkal} and Binder~\cite{binder} evaluate an authorization query
against a fixed trust root and return yes or no. None of them inverts the
question. Given a query that succeeded, none will tell you the minimal set of
principals whose honesty the success depended on. That inversion is the
technical core of this paper, and as far as we can find it has not been done
--- not because it is hard in these systems, but because the systems were
built for a different job. Their trust root is a given; ours is the answer.

\paragraph{Transparency converts trust; it does not remove it.} Certificate
Transparency~\cite{rfc6962,rfc9162} and CONIKS~\cite{coniks} are the canonical
demonstrations that you can take a party you had to trust unconditionally and
turn it into a party whose misbehaviour will be noticed within a bounded time.
This is a genuine and large win, and it is also not the same win as removing
the party. That distinction --- conversion, not removal --- is the one the
tier ladder most needs and most lacks, and it is why detection latency is a
first-class term in our judgment rather than a footnote. A log operator with
a \SI{15}{\minute} gossip bound is still in your trust set. It is in there with
a number next to it.

\paragraph{Hardware roots are assumptions, and the attacks say so
precisely.} Foreshadow~\cite{foreshadow}, SGAxe~\cite{sgaxe},
Plundervolt~\cite{plundervolt} and \AE{}PIC Leak~\cite{aepic}, alongside
Costan and Devadas's architectural analysis~\cite{costan2016}, establish that
enclave isolation is a property that has repeatedly failed. The shape of these
attacks matters more than their existence: they break isolation \emph{without}
invalidating the signing key. The quote remains cryptographically valid while
the claim it makes is false. Any calculus that treats ``the signature
verified'' as discharging the hardware assumption is therefore modelling the
wrong thing, which is why the endorser appears in our output with
$\Delta = \Never$ rather than not appearing at all.

\paragraph{Restructuring trust cryptographically.} Direct Anonymous
Attestation~\cite{daa} is worth naming because it shows how much of a trust
structure can be moved rather than merely documented. We do not do that here;
we compute over the structure as built.

\paragraph{Tabled resolution, which we ended up not needing.} SLG
resolution~\cite{slg} and XSB~\cite{xsb} are the standard answer to a logic
program that diverges on cyclic data, and the design this paper implements
originally specified them for exactly that reason: real deployments contain
mutual and circular attestation, and plain SLD resolution does not terminate
on them. Section~\ref{sec:termination} explains why we did not need it.

% ================================================================
\section{Writing It Down Precisely}
\label{sec:formal}

Go back to the argument in Section~\ref{sec:intro} that does not end. The
reviewer says you are still trusting the build pipeline; the vendor says the
attestation covers it. Both are sincere, and neither has said what ``you are
still trusting the build pipeline'' would have to mean before anybody could
check whether it is true.

So try to write the sentence down and see what it forces on you. Not symbols
for their own sake --- take the thing you actually want to say about a
deployment, ask what it needs in order to be sayable at all, and keep what
survives. Everything below is what the code does: the composition rules are
transcribed from \texttt{src/mechanism.rs} and \texttt{src/solve.rs} in the
artifact, not drafted for the paper.

\subsection{The judgment}

Here is the sentence. \emph{This party has to be honest about this particular
thing, and if they were not, you would find out within this long.} What does
that have to be made of?

It needs a name --- which party. It needs the thing they are being honest
about, kept apart from the name, and it is worth seeing why they have to be
kept apart. ``Intel is trustworthy'' is not a proposition anyone can check.
Intel is being trusted to make silicon that behaves as documented; it is not
being trusted to publish your reference values. Run the two together and you
will not notice the day a party who was fine for one job starts doing
another. So: let $\mathcal{P}$ be a finite set of named principals and
$\mathcal{C}$ a finite set of capability names --- ``signs quotes honestly,''
``publishes correct reference values,'' and so on.

Now the third part, and this is where the shape of everything else gets
decided. How long until you find out? Ask it of the quoting enclave. Suppose
its signing key leaks tonight. How long before somebody notices? Not a long
time --- never. The quotes it forges verify, and they verify tomorrow, and
they verify next year, and nothing anywhere in the system ever produces the
contradiction that would give it away. Write that down as some enormous
number of seconds and you have told a specific lie: that patience is
eventually rewarded. It is not. So a \emph{detection latency} is either a
finite bound or nothing at all:
\begin{equation}
\Delta \in \mathcal{L} = \{\,\mathrm{Bounded}(t) : t \in \mathbb{N}\,\}
                          \cup \{\Never\}.
\end{equation}
$\Never$ is not a large number. It means there is no detection mechanism, so
the violation is silent forever. Keeping it as a distinct symbol rather than
a big integer is deliberate: it is the difference this paper exists to make
visible.

Is a name, a capability and a latency enough? Write a few of them out side by
side and you find two things missing. First: when this party lies,
what does the lie buy? Forging a claim outright is not the same failure as
withholding evidence, and a reader who cannot tell them apart cannot tell
which row of the table to be frightened of. Second, and easier to overlook:
why is this assumption here at all --- which part of the system dragged it
in? Leave that out and you have a list of parties with no way to ask what
would happen if a layer were removed, or whether two layers are leaning on
the same shoulder. So an \emph{assumption} is a four-tuple plus a provenance
tag,
\begin{equation}
a = (p,\ c,\ \Delta,\ \iota,\ m)
  \in \mathcal{P} \times \mathcal{C} \times \mathcal{L}
      \times \mathcal{I} \times \mathcal{M},
\end{equation}
read as: \emph{principal $p$ must be honest about capability $c$; a violation
would be detected within $\Delta$; a breach buys the adversary $\iota$; and
this assumption entered because of mechanism $m$.} The impact $\iota$ ranges
over \textsc{Soundness} (a false claim can be made to verify),
\textsc{Revocation} (revoked or vulnerable components pass as valid),
\textsc{SplitView} (equivocation is possible, but leaves contradictory
evidence), and \textsc{Availability} (evidence can be withheld, not forged).
A \emph{residual trust set} $T$ is a finite set of assumptions.

That is one assumption. What we wanted was a statement about a whole
deployment --- this claim holds, and here is everybody you are still relying
on for it --- which is the same shape with a set in the middle:
\begin{equation}
\Gamma \vdash P : T,\ \Delta
\end{equation}
--- \emph{proposition $P$ holds provided every party in $T$ is honest about
the capability named, and any violation is detected within $\Delta$}. In the
implementation $\Gamma$ is the deployment description: the declared
principals, the mechanisms and their parameters, and the delegation edges. $P$
is the deployment's single \texttt{claim} string. $\Delta$ is derived from
$T$, not carried alongside it: $\Delta = \bigsqcup_{a \in T} \Delta_a$.

\subsection{The two lattices}

A deployment usually has more than one mechanism: a TEE, and a log, and an
anchor over the log, each supporting the same claim and each with its own
answer to the question above. So you are holding several answers and you need
one. What do you do with them?

\paragraph{Latency.} Detection first, because here the wrong answer is the
tempting one. Suppose one layer would report misbehaviour within the hour and
another would never report it at all. What is the detection latency of the
two together? The generous reading says the good layer helps --- there is a
mechanism now, and something is better than nothing. But ask the question the
way an adversary asks it. He is under no obligation to attack the layer that
is watched. He attacks the one that is not, nothing reports anything ever,
and the hourly report goes on arriving punctually about the wrong thing. So
join has to take the worse case and $\Never$ has to absorb rather than
average out. Order $\mathcal{L}$ with every bound below $\Never$ and finite
bounds ordered by magnitude:
\begin{equation}
\Delta_1 \sqcup \Delta_2 =
\begin{cases}
\Never & \text{if } \Delta_1 = \Never \text{ or } \Delta_2 = \Never,\\
\mathrm{Bounded}(\max(t_1, t_2)) & \text{otherwise.}
\end{cases}
\end{equation}
Meet takes the better case, with $\Never$ as the identity.

$\mathcal{L}$ has \emph{infinite} height: $\mathrm{Bounded}(0) <
\mathrm{Bounded}(1) < \cdots$ is an infinite ascending chain below $\Never$.
This matters for Section~\ref{sec:termination} and it is not a problem,
because the fixpoint does not iterate here. $\Delta$ is a fold over the
finished trust set, computed once after the fixpoint has converged; the only
lattice the solver iterates in is the trust-set lattice, whose height is
finite. Confusing the two would make the termination argument false, so it is
worth stating which one carries it.

Read back in words, the join says a system is only as detectable as its least
detectable part. That is a strong claim about the world and we are not
certain of it --- a relying party might reasonably care that \emph{most} of
its exposure is watched, and the max throws that away. Whether it is the
right reading is an open question and Section~\ref{sec:open} returns to it.
The empty trust set has $\Delta = \mathrm{Bounded}(0)$: nothing to detect.

\paragraph{Trust sets.} Now the sets themselves, where the join is the one
nobody argues with: believing a conclusion that rested on two layers means
believing everything either layer needed, so combining is union. The question
that takes work is the other one. Two deployments, two trust sets --- when is
one of them more verifiable than the other?

There is an answer that does not require inventing anything. One deployment
is better if it asks you to believe strictly less: everything in its trust
set is also in the other's, and the other has more besides. That is
containment. So $(\mathcal{T}, \subseteq)$ is the powerset of the assumptions
derivable from a deployment, ordered by inclusion, with join $=$ union and
meet $=$ intersection.

Now notice what you have just agreed to. Inclusion is a \emph{partial} order,
and that is the load-bearing fact of the paper: two trust sets can be
\emph{incomparable}, each containing an assumption the other lacks. Nothing
in the definition promises that any two deployments can be ranked, and it
would be a mistake to add something that did. When two sets come out
incomparable, neither deployment is ``more verifiable'' than the other, and
no ordinal tier can rank them. Whether that case is rare or usual is an
empirical question, and Section~\ref{sec:eval} goes and looks.

We wrote the order by hand rather than deriving it, and the reason is worth
recording because it is a trap. Rust's derived ordering on an ordered set
type is lexicographic and \emph{total}: under it $\{2\}$ is not $\le
\{1,2\}$, so \texttt{join} could return a value that is not above its own
operands, violating the lattice contract in a way that produces wrong answers
rather than a crash.

\subsection{The composition rules}

That is the frame, and by itself it says nothing about any real system. The
content is in one question, asked once per mechanism: \emph{if I believe a
claim because of this, who am I believing, and would I ever find out I was
wrong?} Write $\mathcal{A}(m)$ for the answer --- the assumptions a verifier
must accept in order to believe a claim that mechanism supports.
Table~\ref{tab:rules} is what came back when we asked it of every mechanism the
standard admits. All latencies below are the mechanism's declared parameters,
parsed as durations; \texttt{never} is accepted and means $\Never$.

\begin{table}[h]
\centering\footnotesize
\caption{The composition rules, as implemented in \texttt{src/mechanism.rs}.
Fourteen rows, and the $\Delta$ column divides them three ways. Ten carry an
unconditional $\Never$: the declaration cannot express a bound at all, so it
is a property of the mechanism rather than a default a deployment can
override. Three carry no $\Never$ at all --- the anchoring interval, the
gossip propagation delay and the \texttt{tee\ub attestation} collateral
refresh are declared durations, and a deployment must supply one. The
remaining row, \texttt{anchoring}'s settlement assumption, is the only one
where $\Never$ \emph{is} an overridable default, applied when a settlement
layer is declared without a finality bound.}
\label{tab:rules}
\setlength{\tabcolsep}{4pt}
\begin{tabular}{@{}p{2.4cm}>{\raggedright\arraybackslash}p{5.0cm}>{\raggedright\arraybackslash}p{3.0cm}p{2.3cm}@{}}
\toprule
\textbf{Mechanism} & \textbf{Assumption (principal, capability)} &
\textbf{$\Delta$} & \textbf{Impact} \\
\midrule
\texttt{signing} & signer, \texttt{key\ub custody} & $\Never$ &
\textsc{Soundness} \\
\addlinespace
\texttt{hash\ub chain} & log operator,
\texttt{append\ub only\ub monotonicity} & $\Never$ & \textsc{Soundness} \\
\addlinespace
\texttt{anchoring} & log operator,
\texttt{append\ub only\ub non\ub equivocation} & anchoring interval &
\textsc{SplitView} \\
 & settlement layer (if declared),
\texttt{consensus\ub execution\ub fidelity} & declared finality, else
$\Never$ & \textsc{Soundness} \\
\addlinespace
\texttt{gossip} & each peer, \texttt{view\ub synchronization} &
propagation delay & \textsc{SplitView} \\
\addlinespace
\texttt{witness\ub quorum} & each witness,
\texttt{non\ub collusion\ub $k$\ub of\ub $n$} & $\Never$ &
\textsc{Soundness} \\
\addlinespace
\texttt{tee\ub attestation} & endorser,
\texttt{silicon\ub and\ub microcode\ub integrity} & $\Never$ &
\textsc{Soundness} \\
 & quoting enclave, \texttt{quote\ub signing\ub honesty} & $\Never$ &
\textsc{Soundness} \\
 & collateral authority,
\texttt{accurate\ub collateral\ub issuance} & collateral refresh &
\textsc{Revocation} \\
 & reference-value provider,
\texttt{golden\ub value\ub correctness} & $\Never$ & \textsc{Soundness} \\
 & host, \texttt{measurement\ub injection\ub resistance} & $\Never$ &
\textsc{Soundness} \\
\addlinespace
\texttt{zk\ub proof} & ceremony, \texttt{toxic\ub waste\ub destruction} &
$\Never$ & \textsc{Soundness} \\
 & compiler, \texttt{sound\ub arithmetization} & $\Never$ &
\textsc{Soundness} \\
 & auditor, \texttt{constraint\ub completeness} & $\Never$ &
\textsc{Soundness} \\
\bottomrule
\end{tabular}
\end{table}

Three of these deserve a sentence each.

\emph{\texttt{tee\_attestation} is Table~\ref{tab:tdx}.} The five rows of the
rule are the five rows of the TDX table, and only the collateral authority
carries a bound. This is the sense in which the tool reproduces the
hand-derived result: the deployment file declares five principal identifiers,
a refresh interval and a mechanism kind, and the five-party set --- with each
party's capability, detection bound and impact --- falls out of the rule.

\emph{\texttt{anchoring} is a trade, not a gain.} Everyone knows what
anchoring buys: it converts the log operator's assumption from unbounded to
bounded, which is the transparency-systems result of
Section~\ref{sec:related} made mechanical, and it is a real win. The question
worth asking is what it cost, and you can read the answer off the rule. If a
settlement layer is declared, the rule \emph{adds} an assumption about that
layer's execution fidelity --- a deep reorganization rolls back published
attestation evidence. So the row you were converting is not the only row that
moved. Anchoring is usually described as pure gain; in this calculus it
visibly is not.

\emph{\texttt{witness\_quorum} names every witness, not $k$ of them.} The
capability is \texttt{non\_collusion\_$k$\_of\_$n$} and it is attributed to
each of the $n$. Two things are going on there and they are worth separating,
because Section~\ref{sec:tiers} turns on the difference.

Naming all $n$ is not the arguable part. Drop any witness from the set and
the threshold changes, so every one of them is load-bearing for the claim,
and a trust set that listed only $k$ of them would have to say which $k$ ---
which the deployment cannot know in advance. All $n$ belong.

The \emph{count} is arguable. What a threshold rule asks a verifier to
accept is one proposition, about the set: that no $k$ of these $n$ collude.
Our assumption tuple is keyed by a principal, so it has no way to say ``one
assumption, about this set of $n$.'' A set-valued assumption has nowhere to
live except unrolled across its members, and so $|T|$ charges a quorum $n$
times for something a verifier accepts once. That is a limitation of the
representation rather than a claim about the world, we do not have a fix for
it, and Section~\ref{sec:tiers} shows what it costs. It is also where
Section~\ref{sec:open}'s principal-independence problem bites hardest.

\subsection{Delegation, and why it is a fixpoint}

Mechanisms name principals directly, and if that were all a deployment
contained you could read the trust set off the file with a highlighter. It is
not all. Deployments delegate: $q$ \textsf{speaks-for} $r$ means that
believing $r$ commits you to trusting $q$. Which means the party you were
relying on can hand your reliance to somebody you have never heard of, who
can hand it on again --- so the set is not something you read off, it is
something you have to close. Four rules close it:

\begin{enumerate}[itemsep=1pt,topsep=2pt]
\item A principal named by a mechanism supporting claim $P$ is load-bearing
for $P$.
\item If $r$ is load-bearing for $P$ and $q$ \textsf{speaks-for} $r$, then $q$
is load-bearing for $P$. (Delegation is transitive. This is the rule that
diverges under naive SLD resolution.)
\item Every load-bearing principal contributes its own assumptions to the
residual set.
\item Every delegate carries one assumption of its own: that its delegation
is honest and unrevoked --- capability \texttt{delegation\ub integrity}, with
$\Delta = \Never$ and \textsc{Soundness} impact --- whether or not it has
mechanisms. Delegating to a party you did not otherwise depend on still puts
that party in your trust set.
\end{enumerate}

Rule 4 fires per delegation \emph{edge}, tagged with the specific $r$ it
delegates through, not once per delegate. That detail was a bug before it was
a design: a flat tag made every delegation assumption in a deployment
byte-identical, so a principal delegating into two different mechanisms
collapsed to one deduplicated assumption and a genuine shared dependency
disappeared from the output.

\subsection{Provenance, and why the comparison is order-independent}
\label{sec:orderindep}

The fifth component of an assumption, the mechanism tag $m$, is the one that
has not yet paid for itself. Here is what it is for. Suppose a party turns up
twice in the same trust set. Two quite different things could have happened:
two independent layers each happened to name that party, which is a finding
about the deployment, or one layer named it twice, which is nothing at all.
Strip the tag and those two cases become the same case --- the assumptions
deduplicate into one, the second name vanishes, and the most interesting
result in Section~\ref{sec:eval}, a TEE and a ZK circuit resting on one
compromise, is invisible. The tag is what keeps them apart.

Because $m$ participates in assumption identity, it also participates in
comparison --- and that is where we got it wrong. Tags were originally
positional, of the form \texttt{kind\#index}, index being the mechanism's
position in the deployment file. The consequence is easy to state and was
embarrassing to find: \textbf{a deployment compared as incomparable to a
reordering of itself.} Two files listing the same mechanisms in a different
order produced different tags, hence different assumption tuples, hence two
trust sets neither of which contained the other.

This matters beyond tidiness. The answer this paper offers to the
completeness problem (Section~\ref{sec:open}) is to have a second author
encode the same deployment independently and diff the two trust sets, treating
any divergence as a finding. That experiment is worthless if the comparison
reports divergence for two encodings that differ only in the order somebody
typed the stanzas --- it would produce a divergence every time and mean
nothing.

Tags are now derived from \emph{canonical content}, not position: a mechanism
renders to a fixed-order string with list-valued fields sorted, and repeated
identical declarations are distinguished by an ordinal among mechanisms
sharing that content rather than by file position. So reordering distinct
mechanisms cannot change any tag, while two genuinely repeated declarations
stay distinct. The comparison is order-independent by construction, and the
property is pinned by tests --- one asserting that a deployment and its
reordering compare \textsc{Equal} with an empty diff, one asserting that tags
are a stable multiset under reordering, and one asserting that an
\texttt{a}, \texttt{b}, \texttt{a} sequence assigns the second \texttt{a}
ordinal 1 by content rather than 2 by position.

Removing positional sensitivity did not remove \emph{lexical} sensitivity,
and we found the second one the same way we found the first. Canonical
rendering initially emitted each duration field as the text the author had
typed, so \texttt{12h} and \texttt{720m} --- the same 43\,200 seconds, the
same $\Delta$ in every derived assumption --- produced different tags, and a
deployment was once again incomparable to a rewrite of itself. The same held
for \texttt{never} against \texttt{Never}, since the duration parser is
case-insensitive and the renderer was not. Duration fields are now rendered
from the \emph{parsed} value rather than the source text, which makes the
rendering fallible and is the right trade: an unparseable duration is an
error, not a mechanism whose identity is its own typo. We record this
because the general lesson is not ``do not index by position'' but the
weaker and more useful one: any part of a declaration that reaches a tag
must reach it through the normalising function that defines what that part
\emph{means}, or the tag encodes spelling instead of content.

Two more constructions are refused rather than answered. A deployment
declaring no mechanisms at all solves to the empty
trust set, which takes $\Delta = \mathrm{Bounded}(0)$ and sits below every
other trust set under inclusion --- so a two-line file would be reported as
the most verifiable system in any comparison it entered. And two
byte-identical mechanism stanzas would, under the ordinal rule above, be
tagged distinctly and then read by the shared-dependency analysis as two
independent layers, which two identical declarations never are. Both are
rejected at validation. Neither is a case where a defensible answer exists
and the tool merely picks the wrong one; in both, answering at all is the
error.

The canonical rendering is unambiguous only because principal identifiers
cannot contain the grammar's delimiters. That is enforced once, when a
deployment is validated, and it is the kind of invariant that a later
convenience change could quietly break; the code says so at the point where
it would be broken.

% ================================================================
\section{Why It Terminates}
\label{sec:termination}

The design this implements proposed encoding the calculus as Datalog with
tabled (SLG) resolution~\cite{slg,xsb}, or a Souffl\'{e}
equivalent~\cite{souffle}, because real deployments contain
circular delegation --- $A$ vouches for $B$ vouches for $C$ vouches for $A$
--- and a top-down proof search will chase that cycle forever. That reasoning
is correct about SLD resolution and unnecessary here, and the difference is
worth being precise about, because ``we used a simpler method'' is a claim a
referee is entitled to check.

Here is the argument.

\begin{proposition}[Termination]
For any finite deployment description, the residual trust set is reached in
finitely many steps, regardless of cycles in the delegation graph.
\end{proposition}

\begin{proof}[Argument]
Fix a deployment with a finite set of declared principals $\mathcal{P}$, a
finite list of mechanisms $M$, and a finite set of delegation edges $E$. Let
$A$ be the set of assumptions that can ever be derived. $A$ is finite: each
mechanism $m \in M$ contributes $|\mathcal{A}(m)|$ assumptions, a constant for
every rule except \texttt{gossip} and \texttt{witness\_quorum}, where it is
the number of declared peers or witnesses, and \texttt{anchoring}, where it is
one or two according to whether a settlement layer is declared; and each
delegation edge contributes exactly one \texttt{delegation\_integrity}
assumption tagged with its target. Each of those is finite, so
$|A| \le \sum_{m \in M} |\mathcal{A}(m)| + |E|$ is finite.

The residual set is the least fixpoint of a monotone operator on
$(\mathcal{P}(A), \subseteq)$. Every rule is monotone --- no rule ever removes
an assumption or a load-bearing fact --- and the lattice has height $|A|$,
which is finite. A monotone operator on a complete lattice has a least
fixpoint~\cite{tarski}, and on a lattice of finite height Kleene iteration
reaches it after at most $|A|$ rounds. Cycles in the \textsf{speaks-for}
relation cannot change this, because rule 2 derives facts in the finite domain
$\mathcal{C} \times \mathcal{P}$ rather than generating new terms: a cycle
simply means the same finite set of facts is re-derivable by more than one
path, and re-deriving a fact already present adds nothing and terminates the
round.

The latency lattice takes no part in this. It has infinite height, and if the
solver iterated in it the argument above would not go through; it does not.
$\Delta$ is folded over the converged trust set afterwards, over finitely many
members, so it terminates for the trivial reason.
\end{proof}

The reason tabling is unnecessary, stated plainly: SLD resolution diverges on
cyclic data because it is a top-down search that generates goals, and nothing
stops it from generating the same goal forever. Tabling fixes that by
remembering goals it has already tried. But we are not doing a top-down proof
search. We evaluate bottom-up over ground facts in a finite domain, so the
thing tabling remembers is the thing bottom-up evaluation computes anyway.
The machinery solves a problem this formulation does not have. In practice we
express the rules in \texttt{ascent}~\cite{ascent}, a Datalog embedded in Rust
with support for lattice-valued relations, which does semi-naive bottom-up
evaluation; the lattice-valued relation is exactly the residual trust set,
with union as its join.

This is not only an argument. A test in the artifact builds a delegation cycle
through 300 principals, spanning two mechanism layers, and solves it: every
principal comes out load-bearing, as the cycle requires, and the solve --- the
clock starts after parsing and validation, so it is the fixpoint that is
timed --- returns in well under a second on a laptop in an unoptimized build.
We report that as evidence of termination, not as a performance
result; we have not characterized how the runtime grows, and
Section~\ref{sec:open} lists the complexity bound as unverified.

% ================================================================
\section{Does It Actually Work?}
\label{sec:eval}

Every table and figure in this section is tool output. \texttt{scripts/%
\allowbreak regen-results.sh} runs \parallax{} over the encoded deployments,
writes each one into \texttt{results/}, and this section reads them in from
there. Continuous integration reruns the script and fails the build if a
committed result disagrees with the code that produced it, so a number
\emph{in a table} here cannot outlive that code.

The prose around the tables is a different matter, and worth being exact
about rather than waving at. Where a sentence below restates a figure from a
table --- how many rows are undetectable, how many pairs are N/A --- that
sentence was typed, and a typed restatement can drift away from the table it
sits next to. Each one that could is pinned by a test in the artifact;
\texttt{paper/README.md} names them. That is a weaker guarantee than the
tables have, and we would rather say so than imply the whole section is
generated.

Five deployments are encoded and shipped with the artifact:
$\Sigma_1$, a software-only host with signed measurement logs; $\Sigma_2$,
Intel TDX with ECDSA quotes and PCS collateral; $\Sigma_3$, a 5-of-7 witness
quorum over a transparency log; $\Sigma_4$, a Groth16 rollup with an on-chain
verifier; and $\Sigma_5$, a hybrid combining a TEE with a ZK proof. A sixth
file ships as well: a variant of $\Sigma_5$ whose two layers use genuinely
separate build pipelines, which exists only as the negative control for
Section~\ref{sec:shared}. All of them were written by us. Nothing here has
been validated against somebody else's production deployment, and these
examples are the entire evidence base for every number in this section.

\subsection{Reproducing the hand-derived TDX set}
\label{sec:tdxrepro}

The first success criterion, from the original design, is that the tool
produces the parties of Table~\ref{tab:tdx} from the $\Sigma_2$ description
alone. Before running anything, do it yourself; it takes a minute and it is
the only way to know what the tool is contributing.

Open \texttt{examples/sigma2-tdx.toml}. It declares \SigmaTwoMechanisms{}
mechanism, of kind \texttt{tee\ub attestation}, and besides the kind that
stanza names six things: Intel
endorses the platform, a quoting enclave signs the report, Intel's
certification service supplies the collateral, somebody publishes the
reference values, a cloud operator runs the host --- and the collateral gets
refreshed every twelve hours.

Now apply the rule. To believe the measurement you must believe the silicon
behaves, the enclave signs honestly, the collateral is authentic, the golden
values are right, and the host did not swap the firmware at boot. That is \SigmaTwoAssumptions{} assumptions, one
per name. The refresh interval is not a party; it is a number.

Then ask of each one: how would I find out if it were false? Intel's silicon
is subverted --- the quote still verifies. The reference values are wrong ---
the measurement still matches. The enclave leaks its key --- everything
checks out. Nothing anywhere contradicts any of it. Only the collateral has a
schedule, and only because somebody has to fetch it again.

Count them up: \SigmaTwoParties{} parties, \SigmaTwoUndetectable{} of them
silent. That is what we wrote on a whiteboard before the tool existed, and it
is
Table~\ref{tab:tdx}. Now run the tool. Table~\ref{tab:sigma2} is what
\parallax{} prints, read into this document from the file the tool wrote:
the same parties, the same capabilities, the same one bounded row.

\begin{table}[h]
\centering\small
\caption{The residual trust set \parallax{} computes for $\Sigma_2$, from
\texttt{parallax solve --format latex}. The
\texttt{tee\ub attestation} rule supplies every capability, every impact, and
\emph{which} row is bounded --- \SigmaTwoUndetectable{} of the
\SigmaTwoParties{} are $\Never$ because the
rule says the endorser, the quoting enclave, the host and the
reference-value provider have no detection mechanism. The deployment file
supplies the \SigmaTwoParties{} identifiers and one number: the collateral
refresh interval, which is the bound in the one row that has one.}
\label{tab:sigma2}
% Generated by `parallax solve --format latex` via scripts/regen-results.sh — do not edit.
% A bare tabular: the caption and float live in paper/main.tex.
\begin{tabular}{@{}>{\raggedright\arraybackslash}p{4.6cm}>{\raggedright\arraybackslash}p{5.3cm}ll@{}}
\toprule
Principal & Capability assumed & Detectable in & Impact \\
\midrule
\texttt{did:web:cloud.example.com} & \texttt{measurement\_injection\_resistance} & $\infty$ & \textsc{soundness} \\
\texttt{did:web:intel.com} & \texttt{silicon\_and\_microcode\_integrity} & $\infty$ & \textsc{soundness} \\
\texttt{did:web:pcs.intel.com} & \texttt{accurate\_collateral\_issuance} & \texttt{43200s} & \textsc{revocation} \\
\texttt{did:web:rvp.example.org} & \texttt{golden\_value\_correctness} & $\infty$ & \textsc{soundness} \\
\texttt{urn:qe:tdx} & \texttt{quote\_signing\_honesty} & $\infty$ & \textsc{soundness} \\
\bottomrule
\end{tabular}

\end{table}

If the derivation above felt like a cheat, it was, and the cheat is the
result. The criterion as originally worded --- that the parties emerge
``without those parties being named in the encoding'' --- is not what
happened. We read the names straight out of the file in the first step. So
does the tool.

They are named individually in \texttt{examples/sigma2-tdx.toml}, as the
endorser, quoting enclave, collateral authority, reference-value provider and
host fields of a single \texttt{tee\ub attestation} stanza. What the
deployment file does not contain
is any statement that those five constitute a trust set, or what each of them
must be trusted \emph{for}. That is what the rule supplies: each party's
capability, its detection bound, its failure impact, and --- the part that is
actually a claim about the deployment rather than about the rule --- that
exactly these five and no others are load-bearing for the claim. A sixth
principal declared in the file but not reachable from any mechanism or
delegation does not appear in the output.

So the honest statement of the result is narrower than the original
criterion: the tool derives the \emph{content} of Table~\ref{tab:tdx} from a
mechanism kind and a set of field values, not the \emph{identities} of the
parties. The identities came from whoever wrote the file, which is exactly the
completeness problem of Section~\ref{sec:open} showing up in the success
criterion.

\subsection{The comparison matrix}

Take the two deployments somebody might actually be choosing between:
$\Sigma_2$, the TDX box, and $\Sigma_4$, the Groth16 rollup with its verifier
on chain. Both attest \texttt{execution\ub valid}. That matters --- it is the
same proposition, so the question ``which of these is more verifiable'' is at
least meaningful --- and Section~\ref{sec:formal} says the answer is
containment: one is better exactly when everything in its trust set is also
in the other's. Two sets, one question. Do it by hand.

$\Sigma_2$ we already have: Intel's silicon behaves, the enclave signs
honestly, the collateral is authentic, the golden values are right, the host
does not inject a measurement. Now $\Sigma_4$. Its file declares
\SigmaFourMechanisms{} mechanisms --- a \texttt{zk\ub proof} naming a
ceremony, a compiler and an
auditor, and an \texttt{anchoring} naming a log operator who posts every ten
minutes to a settlement layer with twelve-second finality. Carry those to
Table~\ref{tab:rules} and the assumptions fall out: the ceremony destroyed
its toxic waste, the compiler arithmetized the circuit soundly, the auditor's
constraints are complete, the log operator does not equivocate --- bounded,
at the posting interval --- and the chain executes consensus faithfully,
bounded at finality. That is \SigmaFourAssumptions{} assumptions.

So: is $\Sigma_2$'s set inside $\Sigma_4$'s? Look for Intel's silicon in that
list. It is not there, and it could not be --- a Groth16 verifier does not
care which CPU produced the proof. Then the other direction, which is the one
people forget to check: is $\Sigma_4$'s set inside $\Sigma_2$'s? Look for the
trusted setup ceremony. A TDX deployment does not have one.

Neither set contains the other. Neither deployment is more verifiable than
the other. There is no answer to the question, and that is not the tool being
coy --- it is what the two designs are. They moved their trust to different
places and there is no vantage point from which one of them is above the
other. Run \texttt{parallax compare} on the pair and it prints
\textsc{Incomparable}, which by now should be an anticlimax.

That was one pair, worked by hand. Table~\ref{tab:matrix} is all of them, and
the first thing to notice about it is how few pairs the question can even be
asked of. It is assembled by \texttt{scripts/\allowbreak regen-results.sh}. The
four \textsc{Incomparable} cells are the tool's own output, from
\texttt{parallax compare}. The eight N/A cells are not: the script reads the
\texttt{claim} field out of each deployment file, sees that the two differ,
and writes the N/A row without invoking the tool at all --- because
\texttt{parallax compare} refuses the pair and exits non-zero. Either way the
file is generated rather than typed, and continuous integration regenerates
it and fails on a diff, so it cannot drift from the code.

\begin{table}[h]
\centering\small
\caption{Every ordered pair among the four non-hybrid deployments
$\Sigma_1$--$\Sigma_4$ under set inclusion. Pairs attesting different
propositions are marked N/A rather than omitted: \parallax{} refuses to rank
verifiability across two different claims, and a truncated matrix would be
indistinguishable from a broken run. Eight of the twelve ordered pairs are
N/A, so the table rests on two unordered comparisons ---
$\Sigma_1$/$\Sigma_3$ and $\Sigma_2$/$\Sigma_4$.}
\label{tab:matrix}
% Generated by scripts/regen-results.sh — do not edit.
% A bare tabular: the caption and float live in paper/main.tex.
\begin{tabular}{@{}ll>{\raggedright\arraybackslash}p{6.6cm}@{}}
\toprule
Deployment A & Deployment B & Relation under set inclusion \\
\midrule
\texttt{sigma1-software} & \texttt{sigma2-tdx} & N/A --- different claims (\texttt{measurement\_valid} vs \texttt{execution\_valid}) \\
\texttt{sigma1-software} & \texttt{sigma3-quorum} & Incomparable \\
\texttt{sigma1-software} & \texttt{sigma4-zk} & N/A --- different claims (\texttt{measurement\_valid} vs \texttt{execution\_valid}) \\
\texttt{sigma2-tdx} & \texttt{sigma1-software} & N/A --- different claims (\texttt{execution\_valid} vs \texttt{measurement\_valid}) \\
\texttt{sigma2-tdx} & \texttt{sigma3-quorum} & N/A --- different claims (\texttt{execution\_valid} vs \texttt{measurement\_valid}) \\
\texttt{sigma2-tdx} & \texttt{sigma4-zk} & Incomparable \\
\texttt{sigma3-quorum} & \texttt{sigma1-software} & Incomparable \\
\texttt{sigma3-quorum} & \texttt{sigma2-tdx} & N/A --- different claims (\texttt{measurement\_valid} vs \texttt{execution\_valid}) \\
\texttt{sigma3-quorum} & \texttt{sigma4-zk} & N/A --- different claims (\texttt{measurement\_valid} vs \texttt{execution\_valid}) \\
\texttt{sigma4-zk} & \texttt{sigma1-software} & N/A --- different claims (\texttt{execution\_valid} vs \texttt{measurement\_valid}) \\
\texttt{sigma4-zk} & \texttt{sigma2-tdx} & Incomparable \\
\texttt{sigma4-zk} & \texttt{sigma3-quorum} & N/A --- different claims (\texttt{execution\_valid} vs \texttt{measurement\_valid}) \\
\bottomrule
\end{tabular}

\end{table}

Every pair the tool will rank at all --- that is, every pair attesting the
same proposition --- comes out \textsc{Incomparable}. If that holds up, it is
the headline: under set inclusion there is no partial order over these
designs, so no ordinal tier ladder can be sound, and conformance should
publish a trust set rather than a tier.

\subsection{The shared dependency}
\label{sec:shared}

The second success criterion is a non-obvious finding, so here is the
deployment first and the finding second.

Suppose you are worried about exactly the thing this paper is about. You have
a TDX box, you have read Table~\ref{tab:tdx}, and you do not want your
execution claim resting on Intel's good behaviour alone. So you put a second
layer underneath it that has nothing to do with the first: a zero-knowledge
proof of the state transition. Different mathematics, different suppliers,
no CPU in it anywhere. If the silicon is subverted, the proof still has to
check out. That is $\Sigma_5$, and it is the shape of design that gets sold
as defence in depth.

Write down who you are trusting, layer by layer, straight off the file. From
the TEE stanza: Intel for the silicon, the quoting enclave for the signature,
Intel's certification service for the collateral, the cloud operator for the
host, and \texttt{did:web:buildco.example}, which is the identifier in the
reference-value field. From the ZK stanza: the ceremony for its toxic waste,
the auditor for constraint completeness, and
\texttt{did:web:buildco.example}, which is the identifier in the compiler
field.

Read the two lists again.

The build pipeline is in both. It supplies the golden values the TDX
measurement is compared against, \emph{and} it compiles the circuit whose
constraints the proof relies on. Compromise it once and both layers certify
the same lie --- and the second layer was the entire point of the design.
That is \SigmaFiveAssumptions{} assumptions over \SigmaFiveParties{}
parties, and the gap between those two numbers is where the finding hides.

Nothing in the file remarks on this. Both stanzas name the party in the
ordinary way, in fields that mean unrelated things, and somebody reading top
to bottom has no particular reason to notice that one string occurs twice.

Figure~\ref{fig:shared} is the run, verbatim. The block at the bottom is the
finding: one principal, reached under two different mechanism tags, with the
capability each layer needs from it. The long parenthesised strings are the
canonical mechanism tags of Section~\ref{sec:orderindep} --- they are what
makes the two layers distinguishable, and they are why the same party
appearing twice does not deduplicate into one assumption.

\begin{figure}[h]
\lstinputlisting[basicstyle=\ttfamily\scriptsize, frame=single,
                 framerule=0.3pt, xleftmargin=0.6em, breaklines=true,
                 columns=fullflexible]%
                {../results/hybrid-shared-dependencies.txt}
\caption{\parallax{} \texttt{solve sigma5-hybrid.toml --shared}, as the tool
prints it. The build pipeline is named once by the TEE stanza, as the
reference-value provider, and once by the ZK stanza, as the circuit compiler.
Nothing in the deployment file remarks on that.}
\label{fig:shared}
\end{figure}

We should be careful about how much credit this deserves. We constructed
$\Sigma_5$, so the tool found something we put there. The claim is that the
analysis finds it from the description --- the shared principal is not
remarked on anywhere in the file --- and that a matched variant with genuinely
separate pipelines reports nothing, so the check is not simply always firing.
Whether it finds such things in deployments nobody constructed is untested.

\subsection{Detection latency and policy}
\label{sec:latencyeval}

Of everything you could ask about a deployment, the number a relying party
would actually want is the composed $\Delta$: how long until I would know.
Guess it for $\Sigma_3$ before looking anything up. $\Sigma_3$ is the witness
quorum --- five of seven witnesses over a transparency log, the log anchored
every fifteen minutes, gossip propagating between peers in fifteen seconds.
It is the one design in this set built specifically to make misbehaviour
visible. So: minutes. A quarter of an hour, say.

Now compose it the way Section~\ref{sec:formal} said we must. $\Delta$ is a
join over the whole trust set and the join takes the worse case. The
anchoring interval is a real bound on a real party; so is the gossip delay.
But underneath them sit the seven witnesses, whose capability is that no five
of them collude, and nothing in the system ever reports collusion. One member
at $\Never$ takes the whole system to $\Never$. The fifteen minutes are still
true. They simply do not survive the composition.

Table~\ref{tab:summary} is all five deployments at a glance: how many
assumptions each rests on, how many distinct parties those assumptions name,
how many of them nothing would ever report a violation of, and the composed
$\Delta$.

\begin{table}[h]
\centering\small
\caption{Every encoded deployment, from \texttt{parallax tiers --format
tex-summary}. \emph{Undetectable} counts assumptions with $\Delta = \Never$.
Principals is lower than $|T|$ only for $\Sigma_5$, where one party is
trusted for two different capabilities --- which is the shared dependency of
Section~\ref{sec:shared} showing up as a number.}
\label{tab:summary}
% Generated by `parallax tiers --format tex-summary` via scripts/regen-results.sh — do not edit.
% A bare tabular: the caption and float live in paper/main.tex.
\begin{tabular}{@{}llrrrl@{}}
\toprule
Deployment & Claim & $|T|$ & Principals & Undetectable & Composed $\Delta$ \\
\midrule
\texttt{sigma1-software} & \texttt{measurement\_valid} & 3 & 3 & 3 & $\infty$ \\
\texttt{sigma2-tdx} & \texttt{execution\_valid} & 5 & 5 & 4 & $\infty$ \\
\texttt{sigma3-quorum} & \texttt{measurement\_valid} & 11 & 11 & 7 & $\infty$ \\
\texttt{sigma4-zk} & \texttt{execution\_valid} & 5 & 5 & 3 & $\infty$ \\
\texttt{sigma5-hybrid} & \texttt{execution\_valid} & 8 & 7 & 7 & $\infty$ \\
\bottomrule
\end{tabular}

\end{table}

The last column is the same for every row, and what happened to $\Sigma_3$
happened to all of them: the \emph{Undetectable} column is nonzero in every
row, and one nonzero entry is all it takes. That is the finding rather than a
defect in the table. Whether the max is the right composition is
Section~\ref{sec:open}'s question; what the column shows is that under the
composition we implemented, no deployment we encoded reaches a finite bound.

Now sit in the relying party's chair, because that is who the whole exercise
is for. You are going to write a policy --- a rule you can run in CI before
you accept anybody's attestation --- and you get to say what you will not
carry. What is the first clause? Almost everyone writes the same one: I will
not accept a failure mode that is silent. In this notation, no assumption
with $\Delta = \Never$. Then a second clause so the first is not doing all
the work: whatever is left has to be detectable inside a day.

Those two clauses are \texttt{examples/policy-strict.toml}. Run it against the
five manifests and you get Table~\ref{tab:policy}.

Every one of them fails, and it is worth seeing which clause did it. The
violation counts are the \emph{Undetectable} column of
Table~\ref{tab:summary} over again, entry for entry, so the second clause
never rejected anything: the 24-hour bound does not bind, because the only
assumptions with finite bounds at all are well inside a day. The first clause
did all of it.

\begin{table}[h]
\centering\small
\caption{Each manifest against \texttt{examples/policy-strict.toml}
(\texttt{forbid\ub undetectable = true}, \texttt{max\ub detection\ub latency
= "24h"}), from \texttt{parallax check}. \parallax{} exits 1 on violations,
so this is a gate a relying party could run in CI. On this evidence it is a
gate that rejects everything.}
\label{tab:policy}
% Generated by scripts/regen-results.sh — do not edit.
% A bare tabular: the caption and float live in paper/main.tex.
\begin{tabular}{@{}llr@{}}
\toprule
Deployment & Verdict & Violations \\
\midrule
\texttt{sigma1-software} & violates the policy & 3 \\
\texttt{sigma2-tdx} & violates the policy & 4 \\
\texttt{sigma3-quorum} & violates the policy & 7 \\
\texttt{sigma4-zk} & violates the policy & 3 \\
\texttt{sigma5-hybrid} & violates the policy & 7 \\
\bottomrule
\end{tabular}

\end{table}

A policy that rejects every candidate is not yet a useful policy, and we do
not read this as showing that strict policies are wrong. It shows that the
strict clause everybody reaches for first --- ``no silent failure modes''
--- is not satisfiable by any of the designs we know how to encode, which is
the same fact as Table~\ref{tab:tdx}'s fourth column stated as a
machine-checkable rule. What a relying party would actually write is a
policy that names which undetectable assumptions it will carry and why. The
tool has no way to express that today, and Section~\ref{sec:open} lists it
as open.

% ================================================================
\section{Can the Ladder Be Repaired?}
\label{sec:tiers}

A tier ladder is an ordinal claim. Tier~3 sits above Tier~2, which means
there has to be some way of taking two deployments and saying which is more
verifiable. We can think of four candidates, and they are ours: the list
comes from the design study this paper implements, not from a survey of the
literature, and we do not claim it is exhaustive. If there is a fifth that
works, this section does not rule it out.
\texttt{parallax tiers} runs each of the four over $\Sigma_1$--$\Sigma_4$ ---
the same four the comparison matrix covers --- and reports what each does and
whether the result is an order you could put numbers on.
Table~\ref{tab:orderings} is that output.

None of the four survives. The reason the experiment was worth running is
that they do not fail the same way.

\begin{table}[h]
\centering\small
\caption{The four candidate orderings applied to $\Sigma_1$--$\Sigma_4$,
from \texttt{parallax tiers --format tex-orderings}. The middle column reads
left to right as most verifiable to least, under the candidate's own key.
$\Sigma_5$ is excluded because it is a hybrid built to carry a shared
dependency rather than a fifth design competing for a rung; the comparison
matrix leaves it out for the same reason.}
\label{tab:orderings}
% Generated by `parallax tiers --format tex-orderings` via scripts/regen-results.sh — do not edit.
% A bare tabular: the caption and float live in paper/main.tex.
\begin{tabular}{@{}l>{\raggedright\arraybackslash}p{4.7cm}>{\raggedright\arraybackslash}p{5.6cm}@{}}
\toprule
Candidate ordering & Ranking it produces, most verifiable first & Verdict \\
\midrule
Cardinality $|T|$ & \texttt{sigma1-software} (3) $<$ \texttt{sigma2-tdx} (5) $=$ \texttt{sigma4-zk} (5) $<$ \texttt{sigma3-quorum} (11) & No total order: sigma2-tdx and sigma4-zk tie at 5. \\
\addlinespace
Set inclusion & --- & No total order: of 6 unordered pairs, 4 were refused for attesting different claims and 2 came out incomparable; 0 were strictly ranked and 0 were equal. \\
\addlinespace
Composed $\Delta$ & \texttt{sigma1-software} (never) $=$ \texttt{sigma2-tdx} (never) $=$ \texttt{sigma3-quorum} (never) $=$ \texttt{sigma4-zk} (never) & No order: every deployment scores never, so this key separates nothing. \\
\addlinespace
Collusion cost & --- & Not computable: a deployment description names principals but not what it costs an adversary to corrupt them jointly, which is a property of the adversary and the world rather than of the description. \\
\bottomrule
\end{tabular}

\end{table}

\paragraph{Cardinality ranks the systems backwards.} Count the assumptions
and put the smallest number first. On these four that puts $\Sigma_1$ at the
top: a software-only host that signs its own measurement log and hash-chains
it, with no hardware root, no witnesses, and nothing anybody outside the
host can check. It puts $\Sigma_3$ at the bottom:
a 5-of-7 witness quorum over a log that is anchored and gossiped. Nobody
using the word ``verifiable'' means that. The quorum is doing more work to
be checkable and the count charges it for every witness it added.

The problem is not that the metric is crude. It is that it rewards the wrong
thing. $|T|$ measures how much of a trust structure somebody wrote down, and
under this ordering the way to climb is to write down less --- which is the
completeness problem of Section~\ref{sec:open} turned into an incentive.
Cardinality also fails on its own terms before you get to any of that:
$\Sigma_2$ and $\Sigma_4$ carry the same number of assumptions and tie, so
it is not a total order even where it does rank.

\paragraph{Set inclusion ranks almost nothing.} This is the order the tool
already implements, and nothing about it is wrong ---
inclusion is a real order and it says something true whenever it applies. It
just hardly ever applies. Most of the unordered pairs among these four are
refused outright, because the two sides attest different propositions and
asking which of two different claims is more verifiable is not a question.
Every pair that survives that filter comes out incomparable: each trust set
holds an assumption the other does not. Nothing at all gets ranked. That is
Table~\ref{tab:matrix} restated as an ordering question, and the answer is
that on this evidence the partial order is entirely partial.

\paragraph{Composed detection latency ranks everything the same.} Order by
how fast the system would notice. This one is not wrong either; it is empty.
The composed $\Delta$ is $\Never$ for all four, including the anchored,
gossiped quorum, for the reason Section~\ref{sec:latencyeval} gives: $\Delta$
is a join over the whole trust set, and the \emph{undetectable} column of
Table~\ref{tab:summary} is nonzero in every row. One silent assumption takes
the whole system to $\Never$, and none of these designs gets to zero silent
assumptions. A key that is constant across everything you want to rank
separates nothing.

\paragraph{Collusion cost is not in the input.} Order by what it would cost
an adversary to corrupt enough parties at once. This is the one people reach
for when the others fail, and it is the one \parallax{} refuses to score. A
deployment description names principals. It does not say who owns them, what
jurisdictions they sit in, whose budget could reach them, or whether two
identifiers with different names belong to two companies or one --- which is
the principal-independence problem of Section~\ref{sec:open}, and we do not
know how to compute it either. Those are facts about the world and about a
particular adversary, and the same two designs will rank in opposite orders
for a domestic enterprise and for a cross-border protocol. The tool reports
collusion cost as not computable and says why, rather than printing a number
it would have made up. A made-up number inside a trust manifest is the exact
failure this paper exists to object to.

\paragraph{What follows from that.} Four candidates, four independent
failures: one ranks the systems backwards, one ranks almost nothing, one
ranks everything the same, and one is not derivable from the input at all.
We have no repaired ladder to offer in place of the broken one, and having
run this we do not expect one to exist.

What we have instead is what the original design already proposed, and this
is the result that supports it: conformance should publish a trust manifest
--- the parties, the capabilities, the detection bounds, the impacts --- and
let a relying party evaluate it against a local policy, rather than claim a
tier. A tier is a single number that asserts an ordering.
A manifest asserts no ordering, which on this evidence is the honest thing
for it to assert. Section~\ref{sec:latencyeval} shows the policy check
running and shows it rejecting all five deployments, so this is a direction
rather than a finished answer; Section~\ref{sec:open} lists what it leaves
open, including how an unfamiliar manifest gets evaluated at all.

\paragraph{What this rests on.} Four deployments, all of which we wrote.
The experiment is a computation over four trust sets the tool derived from
four files we authored, and it cannot show that no ordering exists. It shows
that these four candidates fail on these four deployments. Two of the four
failures look structural rather than accidental --- cardinality charges a
deployment for every party it names, and collusion cost is absent from the
input no matter what the input is --- and
we would expect those two to reproduce anywhere. The first comes with a
caveat about our encoding that we set out below. The other two are contingent
on the evidence. Set inclusion is empty here because these four trust sets
happen to be pairwise disjoint or claim-mismatched; a fifth deployment
somebody else wrote could produce a genuine subset pair. Composed $\Delta$ is
constant here because no design we encoded eliminates its last silent
assumption; a design that did would separate immediately. Neither of those
experiments has been run.

One of the two we called structural needs a caveat of its own, and it is the
first thing a referee should push on.

Cardinality puts $\Sigma_3$ last only because \texttt{witness\ub quorum}
attributes \texttt{non\ub collusion\ub $k$\ub of\ub $n$} to each of the $n$
witnesses rather than once to the mechanism. Suppose you reject that and
count one assumption per declared \texttt{[[mechanism]]} stanza instead ---
stanzas, not distinct kinds, and the difference is not idle: $\Sigma_1$
declares \SigmaOneMechanisms{} stanzas of only \SigmaOneMechanismKinds{}
kinds. Under stanzas the ranking does not merely lose its shape; it turns
over. $\Sigma_2$ declares \SigmaTwoMechanisms{} stanza, $\Sigma_4$ declares
\SigmaFourMechanisms{}, and $\Sigma_1$ and $\Sigma_3$ declare
\SigmaOneMechanisms{} apiece. The hardware-attested deployment comes first
and the software-only host ties the witness quorum for last --- which is the
opposite of the paragraph above, not a tie-break away from it.

Count kinds instead and the damage is smaller: $\Sigma_3$ is still last and
$\Sigma_1$ sits mid-table, so ``backwards'' fails without inverting. We lead
with the stanza reading because it is the one that costs us most, and we
would rather a referee found us conceding too much than too little. Neither
reading leaves ``ranks the systems backwards'' standing, which is the part
that matters.

There is a narrower objection that does leave $\Sigma_1$ on top: collapse
only the quorum's $n$ non-collusion assumptions and leave every other rule
alone. But nothing in ``one assumption per mechanism'' singles that reading
out, and we should not get to pick the version of the objection that
flatters us.

So how much of ``cardinality ranks the systems backwards'' survives? The part
that does not depend on the convention is about \emph{principals} rather than
assumptions. Every one of $\Sigma_3$'s witnesses is load-bearing on any
reading --- drop one and the threshold changes --- so the principals column
of Table~\ref{tab:summary} is not an encoding artifact, and it ranks
$\Sigma_3$ last too. What \emph{is} an artifact is the assumption count:
Section~\ref{sec:formal} explains that our tuple is keyed by a principal and
therefore cannot express one assumption about a set of $n$, so a threshold
gets charged $n$ times for a proposition a verifier accepts once.

That leaves us saying less than we would like, and we would rather say the
smaller true thing. Cardinality ranks a deployment worse for naming more
parties, and being checkable by strangers generally means naming more
parties; we think that much holds anywhere. Whether it ranks these four
particular deployments backwards depends on a decomposition we chose, and a
reader who chooses differently gets a different --- and, under the obvious
alternative, an even less defensible --- answer. Neither reading rescues the
ladder, which is the section's point, but only one of them is ours to
assert.

% ================================================================
\section{What We Still Don't Know}
\label{sec:open}

\subsection{The completeness problem, which is the real objection}

A solver computes over the dependencies that were \emph{encoded}. It cannot
discover one that was not.

This is the thing to worry about, and it is worse than an ordinary
limitation, because the failure mode is not silence but confidence. A tool
that prints a tidy five-element trust set while a sixth dependency exists
off-model has converted an open question into a wrong answer, and it has done
so in a format designed to be trusted. That is worse than no tool. If you
take one thing from this paper, we would rather it be this sentence than the
incomparability result.

We do not have a solution. We have three partial answers, none satisfying:

\begin{enumerate}[itemsep=2pt,topsep=2pt]
\item \textbf{Independent encoding and diff.} Have a second author encode the
same deployment without seeing the first encoding, and diff the trust sets.
Any divergence is a finding about the calculus rather than a bug to fix
quietly. The tool supports this --- \texttt{parallax diff} exits non-zero
when two trust sets differ, so it can be a CI gate --- and
Section~\ref{sec:orderindep} explains why the comparison had to be made
order-independent before the experiment could mean anything. We have not run
it with a genuinely independent author. Until we do, this is a plan, not a
result. Note also what it can and cannot detect: it finds dependencies one
author saw and the other missed. It is blind to a dependency that neither
saw, which is the case we are most afraid of.
\item \textbf{Adversarial review of the description.} Treat the deployment
file as the artifact under attack rather than the tool. We have not built a
procedure for this.
\item \textbf{Printing the caveat with every output.} A manifest should carry,
in band, the statement that it enumerates modelled dependencies only, so that
the caveat travels with the artifact instead of living in a paper nobody
reads alongside it. The manifests the tool emits today \emph{do not} carry
such a field; this is the weakest of the three answers, the only one that is
free, and we have not yet implemented it.
\end{enumerate}

\subsection{The rules are not proved sound}

Table~\ref{tab:rules} is a set of design decisions, not theorems. Nothing here
proves that if the parties in $T$ are honest then $P$ actually holds, against
an operational model in which a named principal may deviate arbitrarily.
Soundness is where the artifact's value ultimately has to come from, and it is
future work. What we have is a mechanization of rules that agree with a
careful hand analysis on the cases we checked.

\subsection{Principal independence is asserted, never computed}

A quorum rule reduces to a non-collusion predicate over named parties, and
nothing in the calculus prevents five of seven ``independent'' witnesses from
sharing a jurisdiction, a cloud region, a corporate parent, or a hardware
supplier. The tool finds shared \emph{principals} --- that is the $\Sigma_5$
result --- but a shared principal is the easy case, because at least both
layers named the same identifier. Two distinct identifiers that happen to be
subsidiaries of the same company are invisible to it.

Whether independence can be computed from a description at all, or must be
asserted by somebody who knows the world, is open and consequential. Our
suspicion is that it must be asserted, in which case the honest design is for
the manifest to record who asserted it.

\subsection{Whether $\Delta$ should compose as a max}

We compose detection latency as a join that takes the worse bound, with
$\Never$ absorbing. The argument for it is that a system is only as detectable
as its least detectable part, which is the right thing to tell somebody
deciding whether to rely on the evidence. The arguments against it are real:

Two independent monitors watching the same party do not obviously give you the
worse of their two latencies; you might argue for the better, since either one
detecting is enough. A chain of detection steps --- an anchor that must be
published, then gossiped, then noticed --- might argue for a sum. And a
judgment carries a validity window that the composed $\Delta$ does not
represent: what a $\SI{12}{\hour}$ collateral bound means when the quote is
checked \SI{13}{\hour} later is a question we currently answer by degenerating
to the unmonitored case, which is a choice and not a derivation.

Max is defensible and it is what is implemented. We do not claim it is
derived.

\subsection{The evidence base is five deployments, all ours}

Every result in Sections~\ref{sec:eval} and~\ref{sec:tiers} comes from five
deployments we wrote,
plus a sixth file that exists only as a negative control for one of them. The
tool reproduces a trust set we had already derived by hand --- which is
evidence it implements what we intended, and weaker evidence that what we
intended is right --- and finds a shared dependency in a deployment we
constructed to contain one. Neither of those is validation against reality.
The complexity claim from the original design, that computing $T(\Sigma)$ is
polynomial in principals $\times$ credentials, is likewise unverified: we have
a termination argument (Section~\ref{sec:termination}) and one 300-node
example, not a bound.

\subsection{If the ladder goes, what replaces it}

Supposing incomparability holds and conformance publishes a manifest instead
of a tier, three questions open that we do not answer. How does a relying
party with a local policy evaluate an unfamiliar manifest? The tool has a
policy checker, which is a start and not an answer. How are principals named
stably enough to compare manifests written by different people --- our
examples use DIDs, which pushes the problem somewhere else rather than solving
it. And what stops manifest inflation, where a vendor lists many trivial
assumptions to look rigorous? A trust set is only useful if listing more
assumptions is not rewarded.

% ================================================================
\section{Conclusion}

The claim that hardware attestation leaves nobody to trust is not a small
overstatement; it is wrong in a way that changes what a conformance claim
means. Believing one Intel TDX measurement requires accepting
\SigmaTwoParties{} distinct
parties, and for \SigmaTwoUndetectable{} of them there is no mechanism anywhere in the system
that would ever report a violation. That is worth writing down, and the way
to make sure it gets written down is to compute it rather than argue about it.

The calculus in this paper is small: a judgment with a trust set and a
detection bound, a composition rule per mechanism, two lattices, and a least
fixpoint. Its termination on the case everyone worries about --- circular
delegation --- follows from finite lattice height and needs none of the
machinery originally proposed for it. Its output on the deployments we encoded
suggests that residual trust sets are incomparable rather than ordered, which
would mean the ladder is the wrong abstraction and conformance should publish
a set. We then tried the four orderings that might have rescued the ladder,
and none of them works: one ranks the systems backwards, one ranks almost
nothing, one ranks everything the same, and one asks for a fact no deployment
description contains.

And the tool cannot see what nobody encoded. We would rather that be the
sentence people remember than the incomparability result, because a residual
trust set is a claim about what you depend on, and a confident claim of that
kind is exactly the thing this paper began by objecting to.

% ================================================================
\section*{Acknowledgments}

Prepared within the Proof-of-Control Initiative of the Advanced AI Society.
The Initiative's working groups, board and advisory board are credited
separately; contributions to this manuscript are the basis for authorship.

\bibliographystyle{unsrt}
\bibliography{references}

\end{document}
